\vfill\pagebreak
\section{Variable declaration}
\label{sec:types:variable_declaration}

Inside a function, \token{let} declares a variable: an identifier that stands for
a value. Every value has a type, and the type is what tells the program how to
read the bits and what may be done with them.

A useful picture is a \textit{box}. The variable is the label on the box, the
value is what is inside, and the type is the shape of the box -- which is what
keeps a value of the wrong form from being put in it. This chapter stays with the
types whose values fit in a single box, the scalar types:

\vspace{-10pt}%
\begin{itemize}
\setlength\itemsep{0.2mm}
\item integer types,
\item floating-point types,
\item character types,
\item the boolean type.
\end{itemize}\vspace{-10pt}%

Each of the first three comes in several variants, covering different ranges of
values. The sections that follow take them one at a time; for now, one of each is
enough -- \token{i32} for integers, \token{f32} for floating-point numbers,
\token{c8} for characters -- alongside \token{bool}.

The program below declares one variable of each: \token{x}, \token{y}, \token{z}
and \token{w}.

\begin{lstlisting}[style=coloredverbatim, caption=simple variable declaration, label=lst:types:basic_variables]
use std::io;

fn main() {
  let x: i32 = 10;
  let y: f32 = 3.1415;
  let z: c8 = 'c';
  let w: bool = false;

  println("X : ", x);
  println("Y : ", y);
  println("Z : ", z);
  println("W : ", w);
}
\end{lstlisting}

\input{chapters/chapter2/figures/variable_bins}

Entering \token{main} creates the four boxes of
Figure~\ref{fig:types:variable_bins}, each holding its value. Naming a box hands
back what is inside it, which is what the four \token{println} calls on lines~9
to~12 do. The boxes are a picture and nothing more -- no such object exists at
run time -- but it is the right picture for everything in this chapter.

The type sits between \token{:} and \token{=} in the declaration, and it is
optional: the compiler infers the shape of the box from the first value put in
it. What is not optional is that the shape, once fixed, stays fixed; only values
of that shape go in afterwards.

The four boxes above are also sealed. Their contents cannot be replaced, which
makes them constants. \token{mut}, written after \token{let}, produces a box that
can be opened again -- a \textit{mutable} variable. The book calls both kinds
variables, as everyone does, though strictly it is only the mutable ones that
vary.



Replacing the contents of a mutable variable is the job of \token{=}. The
assignment holds only if both sides have the same type.

The two sides have names. The left operand is the \textit{lvalue}, the right one
the \textit{rvalue} -- \textit{left value} and \textit{right value}. In the box
picture the lvalue is the box and the rvalue is what goes into it. In
Listing~\ref{lst:types:value_change} the lvalue is \token{x}, twice over; the
output follows the listing.

\begin{lstlisting}[style=coloredverbatim, caption=Example of mutable variable, label=lst:types:value_change]
use std::io;

fn main() {
  let mut x = 10;

  x = 42;
  println("X : ", x);

  x = 314;
  println("X : ", x);
}
\end{lstlisting}
\vspace{-10pt}%
\begin{lstlisting}[style=bashVerb, label=lst:types:value_change_output]
X : 42
X : 314
\end{lstlisting}

A second picture explains the run: a cursor travelling down the instructions,
starting at \token{main}.

It creates a box for \token{x} holding \token{10}. One line down, the assignment
puts \token{42} in that box instead. One line down again, \token{println} reads
the box and writes what it finds. Lines~9 and~10 repeat the pair.

That is all imperative programming is -- state that instructions change, one
after another, as control moves through them -- and procedural programming is the
name for doing it with functions. Ymir offers several ways to organise a program,
and all of them rest on this.
